\subsection{The ``Miller-Pera'' emissions dataset}\label{sec:ff_miller}

\begin{description}
%\itemsep1pt\parskip0pt\parsep0pt
%\item \textbf{Global and National Totals}
\item[Global and National Totals]  The Miller-Pera fossil fuel emission
  inventory \citep{pera26a} is derived from independent global total and
  spatially-resolved inventories.  Annual global total fossil fuel
  $\cotwo$ emissions are based on the Appalachian Energy Center’s
  “CDIAC at AppState” project
  (\href{https://rieee.appstate.edu/projects-programs/cdiac/}{https://rieee.appstate.edu/projects-programs/cdiac/};
  \citealp{erb26a}),
  which is an effort to update the original annual global and country
  fossil fuel-$\cotwo$ emissions estimates from the DOE’s Carbon Dioxide
  Information and Analysis Center (CDIAC) \citep{boden17a}. The
  CDIAC at AppState emissions estimates used in \ctver\ extend through
  2022.

\item[Spatial Distribution] Miller-Pera fossil-fuel $\cotwo$ fluxes are
  spatially distributed in two steps: First, the coarse-scale country
  totals through 2022 from CDIAC at AppState are mapped onto a
  \href{http://data.giss.nasa.gov/landuse/country.html}{$\onedeg$
    grid} according to the spatial patterns from the EDGAR 2025 GHG
  release \citep{crippa25a}. The spatial pattern varies
  by year up until the end of the EDGAR product in 2024.  After
  this, the 2024 spatial pattern is held constant.
  Note that while EDGAR provides annual emissions estimates at
  $\onedeg$ resolution, their totals do not agree with those from
  CDIAC at AppState. Thus, only the spatial patterns in EDGAR are
  used, and the total emissions are rescaled to CDIAC values.
  The CDIAC country-by-country totals sum to about 95\% of the
  global total emissions; the remaining 5\% is mapped to global
  shipping routes according to EDGAR, which we treat as a proxy for
  bunker fuel emissions.

\item[Temporal Distribution] For North America between 30 and
  60$\degrees$N, the Miller-Pera system imposes a seasonal cycle derived
  from the first and second harmonics \citep{thoning89a} of the
  \citet{blasing04a} analysis for the United States. The Blasing
  analysis has \textasciitilde{}10\% higher emissions in winter than
  in summer.  This scheme defines a fixed fraction of emissions for
  each month, so while the shape of the annual cycle is invariant, the
  amplitude of that cycle scales with the annual total emissions. For
  Eurasia, a set of seasonal emissions factors from EDGAR distributed
  by emissions sector is used to define fossil fuel seasonality. As
  in North America, this seasonality is imposed only from
  30{-}60$\degrees$N. The Eurasian seasonal amplitude is about 25\%,
  significantly larger than that in North America, owing to the
  absence of a secondary summertime maximum due to air conditioning.
  See Figure~\ref{fig:fossil_tser} for the resulting time series
  of fossil fuel emissions. In order to avoid discontinuities in the
  fossil fuel emissions between consecutive years, a spline curve that
  conserves annual totals \citep{rasmussen91a} is fit to the annual
  emissions in each $\onedeg$ grid cell before the seasonal cycle is
  imposed.

  Whereas early CarbonTracker releases used monthly-constant fossil fuel
emissions, starting with CT2015 we introduced the use of temporal
scaling factors to simulate day-of-week and diurnal variability for
those emissions.  These ``Temporal Improvements for Modeling Emissions
by Scaling'' (TIMES) scaling factors, introduced by \citet{nassar13a},
are again applied to the Miller-Pera monthly emissions estimates for
\ctver.  The scaling factors consist of seven day-of-week global
scaling factor maps, and 24 hourly global scaling factor maps to
represent the diurnal cycle.  For use in TM5, the hourly scaling
factors were aggregated to three-hourly factors to accommodate the
time step of the model.


\item[Extrapolation] The full CDIAC at AppState dataset is
  only available through 2022 at time of running. A prior
  estimate of the fluxes through April of 2026 is required to optimize
  fluxes through the end of 2025, to accommodate our 12-week
  assimilation window.

  For 2023 and 2024 the fractional increases of
  per-country, sectoral emissions are taken from The Energy Institute
  (formerly British Petroleum) Statistical Review of World Energy
  \citep{ei25a} (for coal, oil, gas, and flaring) and the
  \href{https://www.usgs.gov/centers/national-minerals-information-center/cement-statistics-and-information}{USGS
    NMIC Cement Mineral Commodity Summaries} (for cement emissions).
  For example, to calculate 2023 coal emissions for France, the
  EI ratio of 2023 to 2022 coal consumption for France is used to
  scale 2022 CDIAC emissions, which are then used to rescale the
  EDGAR spatial patterns. Only the largest producing countries are
  included in the EI and USGS datasets, so regional or world-average
  ratios are used for extrapolation of other countries.

  For 2025 and the beginning of 2026, global fractional increases
  are assumed (+2.5\% per year for oil and gas, +1\% per year for
  coal and flaring).  The final months, February through April 2026,
  are further scaled by year-on-year ratios from the Carbon Monitor
  near-real-time estimates \citep{liu20a}.  No fuel-type data is
  available from Carbon Monitor, so the same year-on-year fractional
  changes are applied to all fuel types.

\end{description}

